\documentclass[11pt]{article}
\usepackage[utf8]{inputenc}
\usepackage[T1]{fontenc}
\usepackage[french]{babel}
\usepackage{geometry}
\geometry{a4paper, left=2.5cm, right=2.5cm, top=2.5cm, bottom=2.5cm}
\usepackage{fancyhdr}
\usepackage{titling}
\usepackage{hyperref}
\usepackage{xcolor}
\usepackage{titlesec}
\usepackage{fontspec}
\usepackage{ragged2e}
\usepackage{graphicx}
\usepackage[indent=0.66cm]{parskip}
\usepackage{setspace}
\usepackage{enumitem}
\usepackage{amsfonts}
\usepackage{float}
\usepackage{framed}

\title{$title$}
\author{$author$}
\date{\today}

% Configuration des en-têtes et des pieds de page pour les pages normales
\pagestyle{fancy}
\fancyhf{}
\rhead{\vspace{0.66cm}$author$}
\setlength{\headsep}{1.33cm}
\lhead{$shorttitle$}
\rfoot{\thepage}

\definecolor{black}{HTML}{000000} % pour le texte général et les paragraphes
\definecolor{titleColor}{HTML}{273C50} % pour les titres

\setcounter{secnumdepth}{0} % pas de numérotation sur les titres

\setmainfont{Open Sans}
\newfontfamily\oswald{Oswald}
\newfontfamily\sans{Open Sans}

% les titres 1
\titleformat{\section}[block]{\oswald\bfseries\color{titleColor}\fontsize{20}{24}\selectfont}{}{0em}{}
\titlespacing*{\section}{0pt}{20pt}{6pt}
%saut de page avant les titres 1
\let\oldsection\section
\renewcommand{\section}{
    \clearpage
    \oldsection
}

% les titres 2
\titleformat{\subsection}[block]{\oswald\bfseries\color{titleColor}\fontsize{16}{24}\selectfont}{}{0em}{}
\titlespacing*{\subsection}{0pt}{18pt}{6pt}

% les titres 3
\titleformat{\subsection}[block]{\oswald\bfseries\color{titleColor}\fontsize{14}{24}\selectfont}{}{0em}{}
\titlespacing*{\subsection}{0pt}{16pt}{4pt}

% interlignes des paragraphes
\setstretch{1.15}

% listes à puce
\providecommand{\tightlist}{%
  \setlength{\itemsep}{0pt}\setlength{\parskip}{0pt}}
%style des liste à puce
\setlist[itemize,1]{label=\bullet}
\setlist[itemize,2]{label=\diamond}
  
% style des liens
\hypersetup{
    colorlinks=true,
    linkcolor=blue,
    urlcolor=blue,
    citecolor=blue,
    linkbordercolor={0 0 1}, % Couleur du cadre autour des liens internes
    pdfborderstyle={/S/U/W 1} % Style du cadre autour des liens
}

% on ajoute [H] à toute les figures
\let\oldfigure\figure
\renewcommand{\figure}{\oldfigure[H]}

\let\oldincludegraphics\includegraphics
\renewcommand{\includegraphics}[2][]{\oldincludegraphics[width=15cm, #1]{#2}}

\begin{document}

% Créez une page de garde personnalisée sans en-têtes ni pieds de page
\begin{titlepage}
    \centering
    \vspace*{2cm}

    % Insérer les deux images côte à côte dans une seule minipage
    \begin{minipage}{0.5\textwidth}
        \centering
        \includegraphics[width=0.6\textwidth]{ressources/logoUM.png} % Chemin de l'image mis à jour
    \end{minipage}%
    \begin{minipage}{0.5\textwidth}
        \centering
        \includegraphics[width=0.8\textwidth]{ressources/logoIUT.png} % Chemin de l'image mis à jour
    \end{minipage}

    \vspace*{4cm}
    \Huge\oswald\fontsize{26}{24}\color{titleColor}\bfseries $title$\par
    \vspace{2cm}
    \vfill
    \raggedleft
    {\Large\sans\bfseries\linespread{1.2}\fontsize{11}{24}\color{black}\selectfont $schoolyear$\par
    BUT Informatique semestre $semester$\par
    $author$\par}
\end{titlepage}

$body$

\end{document}
